\scanpage{5}
\[
\operatorname{réal}(S,A,R)\lhook\joinrel\longrightarrow E
\]
telle que l'image d'un arc soit un segment.

\uline{Prop} Pour qu'un graphe fini admette une réalisation géom. rectiligne, (dans un vectoriel convenable sur $\mathbb{R}$) il f et s qu'il \deleted{\unsure{a|est}} soit strict -- et alors il admet une réalisation rectiligne dans $\mathbb{R}^{3}$. De façon précise, si $S\hookrightarrow\deleted{\unsure{\mathbb{R}^{3}|R^{3}}}E$ \inserted{$(\simeq\mathbb{R}^{3},\ \unsure{v.\,ex.|p.\,ex.|v.\,a.})$} est une immersion telle que
% Pencil drawing in the left margin beside the conditions. The small dark marks are reproduced without assigning them vertex status.
\reversemarginpar
\margin{\figure{\begin{tikzpicture}[x=.014cm,y=.014cm,line cap=round,line join=round]
\draw[line width=.65pt] (0,0) ellipse (4.6 and 3.4);
\draw[line width=.8pt] (3,3)--(78,86);
\fill (76,84)--(75,91)--(80,88)--(83,90)--(83,84)--(79,82)--cycle;
\draw[line width=.35pt] (88,98)--(147,158);
\draw[line width=.8pt] (147,155)--(147,162);
\draw[line width=.8pt] (-.4,-9)--(.4,-6);
\end{tikzpicture}}}
\normalmarginpar
\[
\forall x,y\in S\quad x\ne y,\quad\text{on ait}\quad
\deleted{\unsure{\{|\exists}}
\underbrace{\overline{xy}}_{\substack{\text{segm}\\\text{joignant}\\x\ \text{à}\ y}}\cap S=\{x,y\},
\]
\inserted{(i.e. trois points \inserted{distincts} jamais alignés)}
\inserted{et $\forall x,y,z,t$ distincts, $(\overline{x,y})\cap(\overline{z,t})=\varnothing$}

$\exists 1$ réalis géom. rectiligne $\unsure{\operatorname{réll}|\operatorname{réel}|\operatorname{réal}}(S,A,R)\lhook\joinrel-E$ qui $f$ \unsure{prolonge|pro\(\overset{?}{l}\)onge}\ldots

Graphe connexe $=$ réal. géom \unsure{connex|connexe} $=$ qui n'est pas \deleted{\unsure{réunion}} ``somme'' de deux autres graphes $\ne\varnothing$

\uline{NB} Tout graphe est de façon unique réunion de graphes connexes \inserted{et disjoints}

[ graphe somme : $S\amalg S',\quad A\amalg A',\quad R\amalg R'\ \ldots$ ]

% A curved arrow in the left margin starts beside NB and points to Graphe, not to the circled insertion multi.
\reversemarginpar
\margin{\figure{\begin{tikzpicture}[baseline=-.5ex,x=1cm,y=1cm]
\draw[->] (.48,1.1) .. controls (-.18,.95) and (-.35,.25) .. (.05,.18) .. controls (.25,.12) and (.48,.08) .. (.78,0);
\end{tikzpicture}}}
\normalmarginpar
% The heading and bracketed degree condition share a baseline; the two additions stand above them.
\[
\begin{array}{@{}l@{\qquad}l@{}}
&\inserted{\text{(ou \uline{contour} (combinatoire))}}\\[-.3ex]
\text{Graphe }\overset{\text{\inserted{\tikz[baseline=(m.base)]{\node (m) at (0,0) {multi};\draw (0,0) ellipse (.50cm and .20cm);}}}}{\text{polygonal}}
&\bigl[\quad\forall s\in S,\quad\operatorname{card}R\{s\}=2\quad\bigr]
\end{array}
\]

Dualité des \deleted{graphes polygonaux} contours (combinatoires).

\deleted{\unsure{donq|donc}} les composantes \inserted{\unsure{conn.|connexes}} d'un \deleted{graphe polygonal}\inserted{contour} sont \deleted{polygonales.} des contours.

Classification des \deleted{graphes polygonaux}\inserted{contours} conn. $\ne\varnothing$

\uline{Prop} Un \deleted{graphe polyg.}\inserted{contour} connexe \inserted{$\ne\varnothing$} est connu à isom. près quand on connait $\operatorname{card}S$ ($=\operatorname{card}A$), qui est n'importe quel cardinal \deleted{fini ou infini}
\[
2\le c\le\aleph_{0}
\]
(strict \unsure{sss|ssi} $c\ne2$ i.e. $c\geqslant3$)
% archive sheet number: 4
